\section{Phạm vi hệ thống (Scope)}

Hệ thống IRMS bao gồm các phân hệ chức năng chính sau:

\subsection{Đặt món số và Quản lý thực đơn}
Nhân viên sử dụng máy tính bảng hoặc thiết bị đầu cuối để nhận đơn nhanh chóng và chính xác. Thực đơn có thể được cập nhật theo thời gian thực (ví dụ: đánh dấu món hết). Đơn hàng bao gồm các hướng dẫn đặc biệt, ghi chú dị ứng, lựa chọn combo và tùy chỉnh theo yêu cầu của khách.

\subsection{Hiển thị đơn bếp và Điều phối quy trình}
Đơn hàng được tự động chuyển từ thiết bị đầu cuối sang Hệ thống Hiển thị Bếp (Kitchen Display System - KDS). KDS tổ chức đơn theo thời gian chuẩn bị, loại món và trạm làm việc (nướng, chiên, tráng miệng). Nhân viên bếp có thể cập nhật tiến độ từ khi bắt đầu nấu đến khi sẵn sàng phục vụ.

\subsection{Quản lý bàn và Đặt chỗ}
Quản lý sơ đồ chỗ ngồi số hóa, theo dõi trạng thái bàn (trống, đang sử dụng, đang dọn dẹp) và ước tính thời gian chờ. Module đặt chỗ hỗ trợ đặt lịch, quản lý danh sách chờ và gửi thông báo cho khách hàng.

\subsection{Thanh toán và Hóa đơn}
Tích hợp quy trình tính hóa đơn bao gồm thức ăn, đồ uống, phí dịch vụ và các chương trình khuyến mãi. Hệ thống hỗ trợ đa dạng phương thức thanh toán (tiền mặt, thẻ, ví điện tử) và cung cấp tùy chọn chia hóa đơn linh hoạt.

\subsection{Giám sát tồn kho}
Nhân viên cập nhật việc sử dụng nguyên liệu sau khi chuẩn bị hoặc phục vụ. Hệ thống tự động cảnh báo khi tồn kho dưới ngưỡng đã định và lưu lại lịch sử sử dụng giúp dự đoán nhu cầu tái đặt hàng.

\subsection{Phân tích và Báo cáo}
Cung cấp các báo cáo doanh thu theo giờ cao điểm, món bán chạy và hiệu suất nhân viên. Phân tích vận hành giúp quản lý nhận diện các nút thắt cổ chai trong quy trình phục vụ để có phương án điều chỉnh kịp thời.
